\chapter{Appendix 4: More about SSL and S-HTTP}

\section{Security Assessment of the SSL Protocol}

\subsection{Introduction}

SSL is a channel security protocol. It provides a secure pipe between
two systems which may be used to carry HTTP transactions or any other
application data. It is not coupled to particular electronic commerce
features or applications. Architecturally, SSL may be compared to
versions of SOCKS which perform encryption or to IP security proposals.
SSL includes the facilities for establishing, negotiating, and managing
the security association between the endpoints (including performing
authentication).\\

SSL operates based on longer term security associations oriented around
the generation of a master key between two parties and short lived
sessions with per-session keys derived from the master key. In the
assessment below, distinctions will be made between the services
applied on a per-association basis vs those performed on a per-session
basis.

\subsection{Security Service Requirements}

\latexhtml{\begin{itemize}}{\begin{htmllist}
      \htmlitemmark{RedBall}}

    \latexhtml{\item Transaction confidentiality}{\item[Transaction
    confidentiality]}
  
  SSL provides confidentiality for all data exchanged between the two
  systems. In SSLv2, encryption (i.e. confidentiality) is mandatory for all
  associations. In SSLv3 confidentiality is an optional service.

  \latexhtml{\item Selective field confidentiality}{\item[Selective field
    confidentiality]}
    
  No selective field confidentiality is provided by SSL. When
  confidentiality is provided, all data exchanged between the parties is
  encrypted.

  \latexhtml{\item Privacy/anonymity}{\item[Privacy/anonymity]}
  
  SSL decouples server and client authentication. Server authentication is
  a mandatory part of security association establishment processing.
  Client authentication is an optional service which may be required and
  invoked by certain servers. If client authentication is not performed,
  the client's identity need not be disclosed.

  \latexhtml{\item Server authentication}
  
  Server authentication is performed as part of the establishment of any
  new association. Authentication is based on the validation of the server
  certificate from which the server's public key is used as part of the
  master key generation process.

  \latexhtml{\item Client authentication}{\item[Client authentication]}
  
  Client authentication may optionally be requested by the server.
  Authentication is based on the validation of the client certificate which
  is used as part of checking the signature on the client's response to an
  optional challenge issued by the server.

  \latexhtml{\item Data integrity}{\item[Data integrity]}
  
  SSL performs a keyed MAC calculation for all data exchanged. The MAC is
  based on a shared secret (created during session establishment), the
  data, and a sequence number. The MAC is mandatory for all information
  exchanged.

  \latexhtml{\item Data origin authentication}{\item Data origin
    authentication}
  
  Data origin authentication is provided as a consequence of the
  authentication performed during session establishment and ensures
  continuity during a session lifetime based on the key management (i.e.
  based on possession of the keys derived from the authenticated session
  establishment.).

  \latexhtml{\item Non-repudiation}{\item Non-repudiation}
  
  No non-repudiation services are provided by SSL.

  \latexhtml{\item Message Stream integrity}{\item[Message Stream
    integrity]}
  
  There are two levels of message stream integrity addressed within SSL.
  The MAC calculated for each record within a session includes a sequence
  number and prevents any modification (insertion, deletion, or reordering)
  of the sequence of records. Multiple sessions may be established under a
  single master key created during the authenticated association
  establishment exchanges. SSL prevents replays of entire sessions by
  exchanging per-session values provided by each party for each session and
  incorporating them into the generation of the per-session key.

  \latexhtml{\item Denial of Service}{\item[Denial of Service]}
  
  Denial of service concerns with SSL relate primarily to consumption of
  resources at the server. An active attacker can readily force a reset of
  a security association and force a new association establishment
  exchange. An attacker could also try to overflow a server's session cache
  and force each new session to repeat the complete association
  establishment exchange. There is little that can be done about either of
  these concerns and neither is significant in the context of what an
  active attacker can do to disrupt communication.
\latexhtml{\end{itemize}}{\end{htmllist}}


\subsection{Other Factors}

\latexhtml{\begin{itemize}}{\begin{htmllist}
      \htmlitemmark{RedBall}}

  \latexhtml{\item Proxy Operation}{\item[Proxy Operation]}
  
  There are substantial issues associated with the use of SSL in a
  proxy/firewall environment. SSL is intended to create a secure channel
  between transport peers. HTTP proxies expect to operate on application
  data. Consequently, there are two choices for proxy operation. Either an
  SSL association can be established through the proxy and viewed as an
  opaque hole through the proxy machine or the proxy can establish SSL
  associations with the client and/or server. In either of these scenarios,
  the proxy must be SSL aware.\\

  In the case of allowing an SSL association to be established through the
  firewall, the security services are provided between the end client and
  server. There are two concerns with respect to the proxy itself.  First,
  there is no way for the firewall to determine what information is flowing
  through the SSL pipe. This amounts to a wide open door into the protected
  environment and is a significant issue. The second concern is the ability
  of the HTTP proxy to perform any application based functions (e.g.,
  access control, filtering, caching). If the SSL association is made
  through the proxy, no application based processing can be performed (in
  fact, the proxy has no way of telling that HTTP transactions are being
  exchanged).\\
  
  If the proxy establishes the SSL association(s), then the proxy can
  perform all of its normal processing (SSL becomes a specialized access
  method). However, the security services are terminated at the proxy.
  Consequently, all information is exposed at the proxy and the
  authentication is limited to the identity of the proxy (not the end
  client or server). There is no structured way for an end client or server
  to determine that a given proxy authenticated by SSL is a ``reasonable''
  proxy for the desired end system. The degree to which this is an issue
  depends on the sensitivity of the information and the access control and
  authentication policies of the clients and servers involved. SSL does not
  provide any means for authenticating multiple parties along a transaction
  path.

  \latexhtml{\item Performance impacts}{\item[Performance impacts]}
  
  There are several dimensions to performance including connection
  establishment latency, transaction overhead, and processing complexity.
  For SSL, the transaction overhead is not significant. The primary impact
  is on session establishment. Between 2 and 4 additional round trip times
  (beyond any other establishment) are required depending upon whether an
  SSL association has already been established (i.e. a master key exists)
  and whether client authentication is performed. In general, the
  operations performed by SSL are not expensive in terms of processing
  (symmetric key and hash operations vs. private key encryption
  operations).\\
  
  SSL requires the maintenance of state information (keys, counters,
  identity, and algorithm information) by both the client and server.
  Management of this state information represents an additional requirement
  for the server.

  \latexhtml{\item Granularity of security services provided by
    protocol}{\item[Granularity of security services provided by protocol]}
  
  Security services are negotiated for an entire association and apply to
  the endpoints of the transport connection. In general, this implies a
  binding to an instance of an application process. The identity being
  authenticated will depend upon the naming information in the client or
  server certificate and on local key management issues (outside of the
  scope of the specification).

  \latexhtml{\item Associated key management and
    cryptography}{\item[Associated key management and cryptography]}
  
  In SSLv2 the authenticated part of the key establishment is limited to
  RSA encryption of the secret key element. In v3, support for additional
  approaches including DH key exchange are included.\\
  
  SSL reuses the same IV throughout an association. In general, this is not
  a problem given the frequency of changing session keys and the protection
  provided by the MAC.

  \latexhtml{\item Negotiation of services, mechanisms, and
    formats}{\item[Negotiation of services, mechanisms, and formats]}
  
  SSL negotiates algorithm information as part of association
  establishment. In SSLv2 the information used to negotiate algorithms is
  not authenticated. This allows an active attacker to cause two parties to
  use a particular selection common to both parties rather than having a
  user's own selection logic invoked. So long as both the client and server
  enforce a check that appropriate protection is applied to all information
  exchanged, this is not a major issue. It may result in weaker
  cryptographic protection being applied to an association than ``should''
  have been used. In SSLv3, all negotiation exchanges are authenticated.
  There are very few options in SSLv2, consequently interoperability
  between independent implementations using common algorithm suites is
  likely.

  \latexhtml{\item Ability to convey certificate and authorization
    information}{\item[Ability to convey certificate and authorization
    information]}
  
  SSL allows only the client or server certificate to be sent. No
  certification path, CRL, or authorization data can be exchanged. It is
  important that certificates used with SSL provide identification
  information which can meaningfully be compared with identity information
  used with HTTP.

  \latexhtml{\item Compromise Recovery and Handling}{\item[Compromise
    Recovery and Handling]}
  
  No explicit mechanisms exist to clean up association data once
  established.  This is an issue only if associations (i.e. master keys)
  exist for an extended period of time. Checks of certificate validity are
  performed only during association establishment.

  \latexhtml{\item Susceptibility to active and passive
    attacks}{\item[Susceptibility to active and passive attacks]}
  
  The security services provided by SSL appear to be robust in the presence
  of active or passive attackers to the extent that robust cryptographic
  algorithms are selected.

  \latexhtml{\item Time based issues}{\item[Time based issues]}
  
  SSL does not rely on any clock based mechanisms for security services.

  \latexhtml{\item Longevity of mechanism}{\item[Longevity of mechanism]}
  
  SSL provides flexibility in algorithm selection and key management which
  provide for longevity of protection. There is a limit on the size of
  challenge information exchanged for association establishment, however
  the size seems more than adequate.

  \latexhtml{\item Interoperability considerations}{\item[Interoperability
    considerations]}
  
  Some aspects of the SSL specification could be clearer in terms of the
  processing expected. Improved clarity would ease the development of
  interoperable implementations.\\
  
  SSL's architectural placement make integration essentially transparent to
  applications (except with respect to exchange of security control
  information) and should simplify integration.

  \latexhtml{\item Error handling and recovery}{\item[Error handling and
    recovery]}
  
  While an error reporting facility exists, there is little structure
  provided. Consequently, applications integrating with SSL may see
  slightly different behavior between different implementations with
  respect to error information returned. SSL is specified to abort a
  session in the event of an error.
\latexhtml{\end{itemize}}{\end{htmllist}}


\section{Security Assessment of the SHTTP Protocol}

\subsection{Introduction}

SHTTP is a transaction security protocol designed explicitly to protect
HTTP transactions. It operates as part of a secured application and
encapsulates (possibly recursively) the application data (i.e. the HTTP
transaction) and applies security services. SHTTP does not rely on an
explicit security association but applies services atomically to each
transaction. Architecturally, SHTTP may be compared to secure messaging
protocols and, in fact, supports the use of such formats to provide the
secure encapsulation. SHTTP includes a description of the means whereby
a party determines what services and formats to apply to a transaction.
SHTTP is designed to support the orthogonal application of security
services (i.e. there are no mandatory services and any combination is
allowed).

\subsection{Security Service Requirements}

\latexhtml{\begin{itemize}}{\begin{htmllist}
      \htmlitemmark{RedBall}}

  \latexhtml{\item Transaction confidentiality}{\item[Transaction
    confidentiality]}
  
  SHTTP provides confidentiality for the encapsulated data exchanged
  between parties. Sufficient (S)HTTP header information to allow
  processing of the SHTTP transaction is sent in the clear.

  \latexhtml{\item Selective field confidentiality}{\item[Selective field
    confidentiality]}
  
  While SHTTP supports recursive encapsulation of transaction data, this is
  expected to be applied to complete HTTP (or SHTTP) transactions.
  Consequently, the granularity of protection remains a complete HTTP
  transaction. This facility could be used to support prior generation of
  signed information objects (contained in minimal HTTP response
  transactions) which would then be returned when requested (possibly with
  additional security encapsulation).

  \latexhtml{\item Privacy/anonymity}{\item[Privacy/anonymity]}
  
  The degree to which anonymity is supported depends upon the security
  services selected. If any authentication or non-repudiation services are
  applied, then an identity of the parties will be disclosed.

  \latexhtml{\item Server authentication}{\item[Server authentication]}
  
  Server authentication may be applied as part of any transaction.
  Authentication may be based on a digital signature applied to the
  transaction as well as based on the key management approach applied to
  encryption keys or based on the keyed MAC calculations. Certain of these
  approaches rely on the validation of certificates.

  \latexhtml{\item Client authentication}{\item[Client authentication]}
  
  Client authentication may be applied as part of any transaction.
  Authentication may be based on a digital signature applied to the
  transaction as well as based on the key management approach applied to
  encryption keys or based on the keyed MAC calculations. Certain of these
  approaches rely on the validation of certificates.

  \latexhtml{\item Data integrity}{\item[Data integrity]}
  
  Data integrity may be applied to any transaction and is based on a
  message integrity check or message authentication code included in the
  protected transaction. Integrity may also be provided by a digital
  signature. If the timestamp based MAC is used, then pre-generation of the
  SHTTP transaction is not possible.

  \latexhtml{\item Data origin authentication}{\item[Data origin
    authentication]}
  
  For SHTTP, data origin authentication is equivalent to the server or
  client authentication service since all SHTTP services are applied to an
  individual transaction.

  \latexhtml{\item Non-repudiation}{\item[Non-repudiation]}
  
  SHTTP supports non-repudiation with proof of origin through the
  application of digital signatures to the encapsulated transaction.  SHTTP
  does not directly support non-repudiation with proof of receipt or any
  other explicit proof of delivery/receipt mechanism.

  \latexhtml{\item Message Stream integrity}{\item[Message Stream
    integrity]}
  
  SHTTP itself applies services to individual transactions and so does not
  explicitly address message stream issues. SHTTP provides a facility to
  allow an application to perform weak sequence related checks through the
  inclusion of a timestamp in MAC check values or strong freshness checks
  through a nonce-based challenge response facility.

  \latexhtml{\item Denial of Service}{\item[Denial of Service]}
  
  Denial of service concerns are limited to the destruction of transactions
  in flight and consumption of resources. The ability of an attacker to
  destroy messages on a channel cannot be addressed by an application
  solution. The consumption of resource issue is based on the ability of an
  attacker to overwhelm a server (or client) which transaction requiring
  extensive processing resources.
\latexhtml{\end{itemize}}{\end{htmllist}}


\subsection{Other Factors}

\latexhtml{\begin{itemize}}{\begin{htmllist}
      \htmlitemmark{RedBall}}

  \latexhtml{\item Proxy Operation}{\item[Proxy Operation]}
  
  SHTTP's selective security services and support for multiple
  encapsulations provide some flexibility for accommodating proxy
  environments. Most importantly, SHTTP allows for authentication (or other
  security services) between multiple parties in a transaction path.
  Consequently, a proxy may be authenticated along with the end client or
  server. Such authentication may be performed either through a single
  signature which can be checked by multiple parties or through multiple
  encapsulations.\\
  
  The interaction with proxies will depend on the peer relationships of
  security services, especially with respect to encryption and key
  management and the operation of the proxy itself. The compatibility with
  proxies will depend upon the degree to which the proxy is SHTTP aware and
  the nature of the processing performed at the proxy. If the proxy is not
  SHTTP aware or if the HTTP transaction is encrypted and hidden from the
  proxy, then the functions which can be performed are very limited (e.g.,
  limited caching, no information based processing).\\
  
  SHTTP does provide an ability to decouple the encryption and
  authentication functions and to encapsulate these services in either
  order so that a proxy may be capable of authenticating the source without
  seeing the data or may examine the data while still preserving end-to-end
  authentication.\\
  
  These abilities will support some, but not all, proxy environments.
  There are limitations using SHTTP with non-SHTTP aware proxies or in
  situations where the proxy must access the information within an HTTP
  transaction.

  \latexhtml{\item Performance impacts}{\item[Performance impacts]}
  
  There are several dimensions to performance including connection
  establishment latency, transaction overhead, and processing complexity.
  For SHTTP the establishment latency is minimal with many situations
  requiring no added exchanges. SHTTP does add a possibly significant
  amount of overhead to a transaction depending on the security format and
  services selected. This overhead is apparent in terms of both the
  increased size of transactions (especially requests) and in the time to
  process a transaction. Depending on the processing, a number of expensive
  (i.e. private key) operations may be required to process a request and
  response. In addition, certification path validation (possibly including
  retrieving certificate validation data) is likely to be required for most
  transactions. The significance of the performance impact must be assessed
  for specific environments (formats, algorithms, and policies).

  \latexhtml{\item Granularity of security services provided by
    protocol}{\item[Granularity of security services provided by protocol]}
  
  Security services are applied to individual HTTP transactions and apply
  to the entire encapsulated data item (an HTTP or SHTTP transaction).  The
  identity being used for security services will depend upon the naming
  information in the client or server certificate (or other data) and on
  local key management issues (outside of the scope of the specification).

  \latexhtml{\item Associated key management and
    cryptography}{\item[Associated key management and cryptography]}
  
  SHTTP provides for a large and extensible set of algorithms and key
  management techniques. These algorithms may be grouped together to
  support selectable suites or may be individually selected.

  \latexhtml{\item Negotiation of services, mechanisms, and
    formats}{\item[Negotiation of services, mechanisms, and formats]}
  
  The SHTTP negotiation process is based on each party making known
  required and optional capabilities. These lists may be sent in HTTP
  exchanges and may also be present in HTML anchors (this will determine,
  in part, whether additional protocol exchanges are required). The sender
  selects an appropriate choice from the list based on the sender's own
  capabilities and policies. The information used to negotiate algorithms
  and formats is not authenticated. This allows an active attacker to cause
  two parties to use a particular selection common to both parties rather
  than having a party's own selection logic invoked. So long as both the
  client and server enforce a check that appropriate protection is applied
  to all transactions, this is not a major issue. It may result in weaker
  cryptographic protection being applied to an association than ``should''
  have been used.\\
  
  SHTTP provides a rich set of options in terms of services, formats, and
  algorithms. Consequently, there is a significant possibility of having
  independent, compliant implementations which are not capable of
  interoperating. Definition of minimum essential requirements (MERs) or
  standard profiles would enhance the likelihood of interoperability.

  \latexhtml{\item Ability to convey certificate and authorization
    information}{\item[Ability to convey certificate and authorization
    information]}
  
  The information which may be conveyed depends upon the format/domain
  selected. Certification path information may always be sent while PKCS-7
  encapsulation also supports transfer of CRLs. Authorization data is not
  currently addressed. It is important that certificates used with SHTTP
  provide identification information which can meaningfully be compared
  with identity information used with HTTP.

  \latexhtml{\item Compromise Recovery and Handling}{\item[Compromise
    Recovery and Handling]}
  
  There is no explicit mechanism to deal with compromise recovery for
  outband or shared secret keys managed externally. This should be dealt
  with through whatever mechanism is used to manage such keys.

  \latexhtml{\item Susceptibility to active and passive
    attacks}{\item[Susceptibility to active and passive attacks]}
  
  The security services provided by SHTTP appear to be robust in the
  presence of active or passive attackers to the extent that robust
  cryptographic algorithms are selected.

  \latexhtml{\item Time based issues}{\item[Time based issues]}
  
  Local clocks are used as one basis of ensuring freshness of transactions.
  Based on the reliability and accuracy of local clocks, this is a weak
  mechanism.

  \latexhtml{\item Longevity of mechanism}{\item[Longevity of mechanism]}
  
  There are no constraints in SHTTP which represent concerns about the
  longevity of the protection, given the flexibility in algorithms.

  \latexhtml{\item Interoperability considerations}{\item[Interoperability
    considerations]}
  
  SHTTP provides a very rich set of options for applying security services.
  The specification describes these tools, but is not clear in guiding the
  use of these options. Additional explanation, and especially the
  inclusion of examples, would enhance the ability to implement compatible
  systems.\\
  
  SHTTP does make use of other standard security encapsulations (PEM, PGP,
  PKCS-7) to apply the security services and this should facilitate
  interoperability and integration in client environments. SHTTP's
  architectural placement requires integration with the HTTP processing
  engine; however, the security processing may be performed by an external
  filter (especially given the use of existing security protocol
  encapsulations).

  \latexhtml{\item Error handling and recovery}{\item[Error handling and
    recovery]}
  
  SHTTP relies on error codes and handling sent in HTTP. This allows
  protection of the error information and makes use of defined processing
  rules which provides a good basis for error handling.
\latexhtml{\end{itemize}}{\end{htmllist}}


\section{Conclusion}

The preceding sections have presented a description of requirements and
other factors to be considered in assessing candidate protocols used to
provide security for electronic commerce applications. The protocols
evaluated here, SSL and SHTTP, represent two possible tools which might
be used to build secure applications. These tools may well need to be
combined with other security mechanisms, particularly secure payment
schemes, to meet overall system security requirements. In addition,
operational security elements beyond the scope of the protocol
specification (e.g. key management, user actions, audit, platform
issues) are critical to system security.\\

In general, SSL and SHTTP provide architecturally distinct, and
possibly complementary, methods for supplying many similar security
services. The main differences lie in four areas:
\begin{htmllist}
  \htmlitemmark{RedBall}
  \item[Operation in proxy environments]
  
    SSL mostly requires a tradeoff between potentially compromising the
    firewall by creating secure associations through the proxy and
    terminating the security services at the proxy. SHTTP may also require
    terminating services at the proxy depending on policy but supports
    security services, especially authentication, between multiple parties
    and provides greater flexibility in cases where the proxy is SHTTP
    aware.

  \item[Nature of integration]
    
    SSL is essentially an added layer between the application and the
    transport service and consequently should be simpler to integrate (and
    may support multiple applications). SHTTP is HTTP specific and is
    likely to require integration with the application package. For
    providing secure system solutions, some degree of coupling between the
    security service and the application will be needed regardless of the
    choice of protocol.

  \item[Performance]
    
    SSL incurs a start up latency, especially for new associations and then
    adds little performance penalty. SHTTP incurs little or no start up
    delay but adds more per transaction overhead in terms of both size and
    processing complexity.

  \item[Non-repudiation]
    
    SSL provides no facilities for non-repudiation while SHTTP provides
    non-repudiation with proof of origin when digital signatures are
    employed.
\end{htmllist}

This document serves to provide information to system integrators and
providers on the characteristics of the respective protocols. The
selection of one or more as elements of a system solution will depend
on the requirements of those particular systems and on a detailed
analysis of customer, operational, and regulatory characteristics.

